% Book 4, Year 2, chapter 18 -- The network layer: IP and routing.
% Provenance (never printed): sources/book4/BRIEFS.md, chapter 18.
\chapter{The Network Layer: IP and Routing}\label{ch:b2:network-layer}

A packet leaves the laptop at the library for a server at the far end of town. Each \omterm{def:b2:network-layer:routing}{router} on the way takes one off
a counter in the packet, and the \omterm{def:b2:network-layer:routing}{router} that brings it to zero sends back a note. Sending packets with counters 1,
2, 3, \dots\ therefore reveals the path, one \omterm{def:b2:network-layer:routing}{router} at a time:

\omoutput{code/bachelor-2/18-network-layer/outputs/traceroute_town.txt}{The town's network, traced (a simulation,
with one \omterm{def:b2:network-layer:routing}{router} at each place of the town).}\label{out:b2:network-layer:hook}

Nine \omterm{def:b2:network-layer:routing}{routers}, then the server: the laptop knew none of them, and each \omterm{def:b2:network-layer:routing}{router} chose the next one on its own. The \omterm{def:b2:link-layer:frame}{link layer} of
\cref{ch:b2:link-layer} carries a frame across one link; the \emph{network layer} carries a packet across many,
from any machine to any other. This chapter covers its addresses, how a \omterm{def:b2:network-layer:routing}{router} chooses where a packet goes next, and
how \omterm{def:b2:network-layer:routing}{routers} learn the routes in the first place.

\begin{recall}
Book 2: an IPv4 address is four numbers from 0 to 255; a mask says which part names the network; a \omterm{def:b2:network-layer:routing}{routing} table,
read by hand, tells a \omterm{def:b2:network-layer:routing}{router} where to send a packet, and \omterm{def:b2:network-layer:routing}{routers} learn routes either from their neighbours' distances
or from a map of the whole network. Year 1: Dijkstra's algorithm and the Bellman--Ford algorithm compute shortest
paths. \cref{ch:b2:link-layer}: the \omterm{def:b2:link-layer:inet-checksum}{Internet checksum}; frames, MAC addresses and ARP.
\end{recall}

\section{Addresses and prefixes}\label{sec:b2:network-layer:addresses}

An IPv4 address is a 32-bit number, written as four bytes in decimal: \code{203.0.113.77} is the number whose
bits are
\[ 11001011\;00000000\;01110001\;01001101 . \]
Addresses are handed out in blocks that share their first bits.

\begin{definition}[Prefix]\label{def:b2:network-layer:prefix}
% ledger: b2:cidr-longest-match
A \emph{prefix}\index{prefix (IP)} $a/p$, with $0 \le p \le 32$, is the set of addresses whose first $p$ bits are
those of $a$. Its \emph{mask} is the 32-bit number made of $p$ ones then $32 - p$ zeros. An address $x$ belongs to
$a/p$ exactly when $x \mathbin{\&} \text{mask} = a \mathbin{\&} \text{mask}$; that common value is the
\emph{network address}, and the address with the $32 - p$ last bits all ones is the \emph{broadcast address}.
\end{definition}

\begin{omfigure}
\begin{bytefield}[bitwidth=0.95em]{32}
\bitheader{0,7,8,15,16,23,24,25,26,31} \\
\bitbox{8}{\scriptsize 203} & \bitbox{8}{\scriptsize 0} & \bitbox{8}{\scriptsize 113} & \bitbox{2}{\scriptsize 01}
& \bitbox{6}{\scriptsize 001101} \\
\bitbox{26}{\scriptsize network part: 26 bits, the prefix 203.0.113.64/26} & \bitbox{6}{\scriptsize host: 13}
\end{bytefield}
\par\smallskip{\small The address 203.0.113.77 in the \omterm{def:b2:network-layer:prefix}{prefix} /26: the first 26 bits name the network, the last 6 a
machine in it. The mask is 255.255.255.192.}
\end{omfigure}

\omcode{code/bachelor-2/18-network-layer/listings/cidr.py}{16}{25}{The mask, the network and the broadcast
address}

\omoutput{code/bachelor-2/18-network-layer/outputs/cidr.txt}{Three \omterm{def:b2:network-layer:prefix}{prefixes}, each checked against Python's
\code{ipaddress}.}\label{out:b2:network-layer:cidr}

\begin{method}[Subnet arithmetic]\label{met:b2:network-layer:subnet}
A \omterm{def:b2:network-layer:prefix}{prefix} $/p$ holds $2^{32-p}$ addresses, of which $2^{32-p} - 2$ can name machines (the network and broadcast
addresses are kept). To split $a/p$ into $2^k$ equal subnets, extend the \omterm{def:b2:network-layer:prefix}{prefix} by $k$ bits: the subnets are $a + i
\cdot 2^{32-p-k}$ over $/(p + k)$, for $i = 0, \dots, 2^k - 1$. To find the smallest \omterm{def:b2:network-layer:prefix}{prefix} for $h$ machines, take
the least $s$ with $2^s - 2 \ge h$: the \omterm{def:b2:network-layer:prefix}{prefix} is $/(32 - s)$.
\end{method}

% ledger: b2:doc-ipv4
% ledger: b2:private-blocks
Some blocks are set aside. Every address printed in this chapter comes from the three blocks reserved for
documentation, \code{192.0.2.0/24}, \code{198.51.100.0/24} and \code{203.0.113.0/24}, except for the private
blocks \code{10.0.0.0/8}, \code{172.16.0.0/12} and \code{192.168.0.0/16}, which any home or company may use inside
its own network and which never appear on the Internet itself (\cref{sec:b2:network-layer:nat}).

\begin{history}
% ledger: b2:ipv4-exhaustion
$2^{32}$ is about 4.3 billion addresses, fewer than the machines that want one. On 3 February 2011 the body that
hands out address blocks gave its last five blocks of $2^{24}$ addresses, one to each regional registry, and had none
left. Address translation (\cref{sec:b2:network-layer:nat}) and IPv6 (\cref{sec:b2:network-layer:ipv6}) are the two
answers.
\end{history}

\section{Forwarding}\label{sec:b2:network-layer:forwarding}

\begin{definition}[Router, forwarding, routing]\label{def:b2:network-layer:routing}
A \emph{router}\index{router} is a machine with interfaces on several links that passes packets from one to another. It
does two different jobs. \emph{Forwarding}\index{forwarding} is what it does for each packet, millions of times a
second: look at the destination address and send the packet out of one interface. \emph{Routing}\index{routing} is how
it fills the table it forwards with, by exchanging messages with other routers (\cref{sec:b2:network-layer:routing}).
\end{definition}

\begin{definition}[Forwarding table, longest-prefix match]\label{def:b2:network-layer:lpm}
A \emph{forwarding table}\index{forwarding table} is a list of \omterm{def:b2:network-layer:prefix}{prefixes}, each with a next hop. A packet for address
$x$ goes to the next hop of the \emph{longest} \omterm{def:b2:network-layer:prefix}{prefix} of the table that contains $x$
(\emph{longest-prefix match}\index{longest-prefix match}). The \omterm{def:b2:network-layer:prefix}{prefix} \code{0.0.0.0/0}, which contains every
address, is the \emph{default route}.
\end{definition}

The longest \omterm{def:b2:network-layer:prefix}{prefix} wins because it is the most precise: a provider announces its whole block, and one customer
inside it, reached another way, announces its own smaller block. Trying every \omterm{def:b2:network-layer:prefix}{prefix} costs one test per entry, and
backbone \omterm{def:b2:network-layer:routing}{routers} hold hundreds of thousands of them. A \emph{binary trie} does better: a tree in which each node has
a child for bit 0 and one for bit 1, and each \omterm{def:b2:network-layer:prefix}{prefix} of length $p$ marks the node reached by its $p$ bits.

\omcode{code/bachelor-2/18-network-layer/listings/lpm_trie.py}{5}{24}{\omterm{def:b2:network-layer:lpm}{Longest-prefix match} on a binary
trie}\label{lst:b2:network-layer:trie}

\begin{omfigure}
\begin{tikzpicture}[font=\scriptsize, n/.style={circle, draw, minimum size=6mm, inner sep=0pt},
  h/.style={n, fill=omDef!30}, e/.style={draw, gray}]
\node[h] (r) at (0,0) {0};
\node[n] (a) at (-2,-1.1) {}; \node[h] (b) at (2,-1.1) {1};
\node[h] (a1) at (-1,-2.2) {3}; \node[n] (b0) at (1,-2.2) {};
\node[h] (b01) at (2,-3.3) {2};
\draw[e] (r) -- node[above left] {0} (a); \draw[e] (r) -- node[above right] {1} (b);
\draw[e] (a) -- node[above right] {1} (a1); \draw[e] (b) -- node[above left] {0} (b0);
\draw[e] (b0) -- node[above right] {1} (b01);
\draw[omThm, very thick, -{Stealth[length=2mm]}] (r) -- (b);
\draw[omThm, very thick, -{Stealth[length=2mm]}] (b) -- (b0);
\node[anchor=west, align=left, text=omThm] at (3.0,-2.2) {lookup of 1001\dots: root, 1, 10;\\ no child for 0 at
10;\\ the last mark met is 1: port 1};
\end{tikzpicture}
\par\smallskip{\small The trie of the table $\ast \to 0$, $01\ast \to 3$, $1\ast \to 1$, $101\ast \to 2$ (shaded
nodes carry a next hop). A lookup walks down the bits of the address and remembers the last mark it met.}
\end{omfigure}

\begin{proposition}[Lookup in a trie]\label{prop:b2:network-layer:lpm}
\code{lookup} returns the next hop of the longest \omterm{def:b2:network-layer:prefix}{prefix} of the table that contains the address, after at most 32
steps.
\end{proposition}

\begin{proof}
After $k$ steps the walk is at the node of the address's first $k$ bits. The node of a \omterm{def:b2:network-layer:prefix}{prefix} $a/p$ lies on this
path exactly when the first $p$ bits of the address are those of $a$, that is, when the \omterm{def:b2:network-layer:prefix}{prefix} contains the address.
The walk passes the marked nodes of all these \omterm{def:b2:network-layer:prefix}{prefixes} in order of length, keeping the last, and stops only when the
path leaves the trie, below which no node is marked, or after 32 bits. So it returns the longest. Each step consumes
one bit of a 32-bit address.
\end{proof}

\omoutput{code/bachelor-2/18-network-layer/outputs/lpm_trie.txt}{Five lookups in a table of five
\omterm{def:b2:network-layer:prefix}{prefixes}.}\label{out:b2:network-layer:lpm}

\omoutput{code/bachelor-2/18-network-layer/outputs/lpm_bench.txt}{Lookup time against table size on the laptop,
for 500 random addresses.}\label{out:b2:network-layer:bench}

The linear scan's time grows with the table; the trie's barely moves. Hardware \omterm{def:b2:network-layer:routing}{routers} use compressed tries, and
special memories that compare an address with every \omterm{def:b2:network-layer:prefix}{prefix} at once.

\subsection{The IPv4 header}

\begin{omfigure}
\begin{bytefield}[bitwidth=0.95em]{32}
\bitheader{0,4,8,16,19,31} \\
\bitbox{4}{\scriptsize version} & \bitbox{4}{\scriptsize IHL} & \bitbox{8}{\scriptsize type of service} &
\bitbox{16}{\scriptsize total length} \\
\bitbox{16}{\scriptsize identification} & \bitbox{3}{\scriptsize flags} & \bitbox{13}{\scriptsize fragment offset}
\\
\bitbox{8}{\scriptsize time to live} & \bitbox{8}{\scriptsize protocol} & \bitbox{16}{\scriptsize header checksum}
\\
\bitbox{32}{\scriptsize source address} \\
\bitbox{32}{\scriptsize destination address}
\end{bytefield}
\par\smallskip{\small The IPv4 header without options: 20 bytes, five rows of 32 bits. IHL is the header's length in
rows; the protocol says what the payload is (6 for TCP, 17 for UDP).}
\end{omfigure}

% ledger: b2:ipv4-header
% ledger: b2:ipv4-checksum
% ledger: b2:ipv4-ttl-zero
Every packet starts with this header. Two fields change at every \omterm{def:b2:network-layer:routing}{router}.

\begin{definition}[Time to live]\label{def:b2:network-layer:ttl}
The \emph{time to live}\index{time to live} of a packet is a counter in its header that every \omterm{def:b2:network-layer:routing}{router} decreases by
one; a \omterm{def:b2:network-layer:routing}{router} that brings it to zero destroys the packet, so that a packet caught in a \omterm{def:b2:network-layer:routing}{routing} loop cannot circle
forever, and sends back a ``time exceeded'' message.
\end{definition}

That message is what \cref{out:b2:network-layer:hook} used. The \emph{header \omterm{def:b2:link-layer:edc}{checksum}}, the \omterm{def:b2:link-layer:inet-checksum}{Internet checksum} of \cref{ch:b2:link-layer} computed over the header with the
field taken as 0, must then be computed again.

\omcode{code/bachelor-2/18-network-layer/listings/ipv4_header.py}{26}{33}{What a \omterm{def:b2:network-layer:routing}{router} does to each header}

\omoutput{code/bachelor-2/18-network-layer/outputs/ipv4_header.txt}{A header built, decoded and forwarded once.}

The \omterm{def:b2:network-layer:ttl}{time to live} dropped from 64 (\code{0x40}) to 63, and the \omterm{def:b2:link-layer:edc}{checksum} went up by \code{0x0100}, the same amount in
the same byte. Error reports such as ``time exceeded'' or ``destination unreachable'' travel in small packets of
their own protocol, ICMP, which is also what \code{ping} sends.

\section{Address translation}\label{sec:b2:network-layer:nat}

A home has one public address and many machines with private ones. Its \omterm{def:b2:network-layer:routing}{router} translates.

\begin{definition}[Network address and port translation]\label{def:b2:network-layer:nat}
% ledger: b2:napt
A \emph{NAT}\index{network address translation} \omterm{def:b2:network-layer:routing}{router} rewrites each packet that leaves the private network: the
private source address and port become its own public address and a port it chooses, and it records the pair in a
table. A packet that comes back to that public port is rewritten the other way and delivered. A packet that arrives
for a port with no entry has nobody to go to, and is dropped.
\end{definition}

\omcode{code/bachelor-2/18-network-layer/listings/nat_sim.py}{5}{20}{The translation table}

\omoutput{code/bachelor-2/18-network-layer/outputs/nat_sim.txt}{Two machines inside, one address
outside.}\label{out:b2:network-layer:nat}

\begin{omfigure}
\begin{tikzpicture}[font=\scriptsize, b/.style={draw, rounded corners=2pt, align=center, inner sep=3pt},
  w/.style={-{Stealth[length=1.8mm]}, thick}]
\node[b, fill=omDef!12] (h) at (0,0) {machine inside\\ 192.168.1.11};
\node[b, fill=omThm!12] (n) at (4.6,0) {home router (NAT)\\ outside: 203.0.113.7};
\node[b, fill=omDef!12] (s) at (9.2,0) {server\\ 198.51.100.20};
\draw[w] (h.north east) -- node[above, align=center] {from 192.168.1.11:50000\\ to 198.51.100.20:80}
  (n.north west);
\draw[w] (n.north east) -- node[above, align=center] {from 203.0.113.7:40001\\ to 198.51.100.20:80}
  (s.north west);
\draw[w, omDef] (s.south west) -- node[below, align=center] {to 203.0.113.7:40001} (n.south east);
\draw[w, omDef] (n.south west) -- node[below, align=center] {to 192.168.1.11:50000} (h.south east);
\node[draw, dashed, below=9mm of n, align=center] {table: 40001 $\leftrightarrow$ 192.168.1.11:50000};
\end{tikzpicture}
\par\smallskip{\small One connection through a \omterm{def:b2:network-layer:nat}{NAT}. The server sees only the \omterm{def:b2:network-layer:routing}{router}'s address; the table maps the
\omterm{def:b2:network-layer:routing}{router}'s port back to the machine inside.}
\end{omfigure}

Translation saved the Internet from running out of addresses, at a price. A machine inside cannot be reached until
it has spoken first, so two machines both behind \omterm{def:b2:network-layer:nat}{NATs} cannot simply connect to each other; a protocol that writes
addresses inside its messages sees them unchanged while the headers are rewritten; and the \omterm{def:b2:network-layer:routing}{router} must keep state for
every conversation, which the end-to-end design of the Internet had avoided.

\section{Fragmentation}\label{sec:b2:network-layer:fragmentation}

\begin{definition}[Maximum transmission unit, fragmentation]\label{def:b2:network-layer:fragmentation}
The \emph{maximum transmission unit}\index{MTU} (MTU) of a link is its largest frame payload (1500 bytes on Ethernet).
\emph{Fragmentation}\index{fragmentation} cuts a packet larger than the next link's MTU into \emph{fragments}, each with
its own header, which the destination puts back together.
\end{definition}

An IPv4 \omterm{def:b2:network-layer:routing}{router} fragments a packet whose ``don't fragment'' flag is not set.

% ledger: b2:ipv4-frag-offset
% ledger: b2:ipv4-576
\omoutput{code/bachelor-2/18-network-layer/outputs/fragment.txt}{4000 bytes of payload through links of \omterm{def:b2:network-layer:fragmentation}{MTU} 1500 and
576.}\label{out:b2:network-layer:fragment}

Each fragment keeps the packet's identification; its \emph{offset} says where its data starts, in units of 8 bytes,
which is why every fragment but the last carries a multiple of 8 bytes; the \emph{more fragments} flag is 1 on all
but the last. One lost fragment loses the whole packet. Modern hosts avoid \omterm{def:b2:network-layer:fragmentation}{fragmentation}: they send with ``don't
fragment'' set and lower their packet size when a \omterm{def:b2:network-layer:routing}{router} reports that a packet was too big. Every host must accept
packets of up to 576 bytes.

\section{IPv6}\label{sec:b2:network-layer:ipv6}

% ledger: b2:ipv6-header
\begin{definition}[IPv6 address]\label{def:b2:network-layer:ipv6}
An \emph{IPv6 address}\index{IPv6 address} is an address of 128 bits, written in the text form given below.
\end{definition}

$2^{128}$ addresses are enough to give every machine a public one. The header is simpler
than IPv4's: a fixed 40 bytes with version, traffic class, flow label, payload length, next header (the protocol
field), hop limit (the \omterm{def:b2:network-layer:ttl}{time to live}), and the two addresses. There is no header \omterm{def:b2:link-layer:edc}{checksum}, since the link and
transport layers each have their own, and \omterm{def:b2:network-layer:routing}{routers} never fragment: only the sender does, and every link must carry
at least 1280 bytes.

% ledger: b2:ipv6-text
% ledger: b2:doc-ipv6
An address is written as eight groups of 16 bits in hexadecimal, separated by colons. Leading zeros in a group are
dropped, and one run of at least two groups of zeros, the longest (the first if two are equally long), is written
\code{::}. Letters are in lowercase. The \omterm{def:b2:network-layer:prefix}{prefix} \code{2001:db8::/32} is reserved for documentation.

\omoutput{code/bachelor-2/18-network-layer/outputs/ipv6_text.txt}{Short forms expanded to eight groups, as
\code{ipaddress} does.}

\section{Routing}\label{sec:b2:network-layer:routing}

\omterm{def:b2:network-layer:lpm}{Forwarding tables} are filled by \omterm{def:b2:network-layer:routing}{routing} protocols. Inside one network, \omterm{def:b2:network-layer:routing}{routers} use one of two families, both
computing shortest paths in the graph of \omterm{def:b2:network-layer:routing}{routers} and links, where each link has a cost.

\subsection{Distance vector}

\begin{definition}[Distance-vector routing]\label{def:b2:network-layer:dv}
In \emph{distance-vector routing}\index{distance vector routing}, each \omterm{def:b2:network-layer:routing}{router} knows only its neighbours and the cost of
its links to them, and repeatedly sends its table of distances to its neighbours, keeping for each destination the
neighbour that offers the shortest total. A \emph{count to infinity} is the slow rise of distances after a failure;
\emph{split horizon} is a remedy, below.
\end{definition}

 Periodically, every \omterm{def:b2:network-layer:routing}{router} sends its table of distances to its neighbours, and rebuilds its own
table from theirs: the distance to $d$ is the least, over neighbours $v$, of the link cost to $v$ plus $v$'s distance
to $d$, and the next hop is that $v$. This is the Bellman--Ford algorithm, run by all the \omterm{def:b2:network-layer:routing}{routers} together.

\omcode{code/bachelor-2/18-network-layer/listings/distance_vector.py}{5}{18}{One synchronous round of distance
vector}

\omoutput[1-5]{code/bachelor-2/18-network-layer/outputs/distance_vector.txt}{Four \omterm{def:b2:network-layer:routing}{routers}, links A--B 1, B--C 2,
A--C 4, C--D 1; the tables after each round.}\label{out:b2:network-layer:dv}

After round 1 each \omterm{def:b2:network-layer:routing}{router} knows its neighbours; A reaches C at cost 4 directly. In round 2 it hears B's distance 2
to C and switches to the path through B, cost 3. After round 3 nothing changes.

\begin{theorem}[Convergence]\label{thm:b2:network-layer:dv-converge}
Let the network have $n$ \omterm{def:b2:network-layer:routing}{routers} and positive link costs, and let every table start with only the \omterm{def:b2:network-layer:routing}{router} itself, at
distance 0. After round $k$ of synchronous distance vector, each \omterm{def:b2:network-layer:routing}{router}'s distance to each destination is the least
cost of a path to it with at most $k$ links. Hence after $n - 1$ rounds every distance is the shortest-path distance,
and no table changes after that.
\end{theorem}

The proof, by induction on $k$, is \cref{exo:b2:network-layer:8}. On the town's 30 \omterm{def:b2:network-layer:routing}{routers} the tables stop changing
after 8 rounds, the most links on any shortest path, and agree with Dijkstra's algorithm run at each \omterm{def:b2:network-layer:routing}{router}:

\omoutput[1-3]{code/bachelor-2/18-network-layer/outputs/town_network.txt}{Distance vector against link state on the
town's network.}

Good news travels fast. Bad news does not. Take a separate line A--B--C, both links of cost 1. When the link from B to C fails, B no longer reaches C
directly, but A still announces C at distance 2, through B. B takes that route, at distance 3; A then hears 3 from
B, and so on:

\omoutput[1-1,6-21]{code/bachelor-2/18-network-layer/outputs/distance_vector.txt}{The count to infinity.}

\begin{proposition}[Count to infinity]\label{prop:b2:network-layer:count-to-infinity}
In this situation, where A's only route to C goes through B and B's through A, the larger of their two distances to C
grows by exactly the cost of the link A--B at each round, until it reaches the value taken as infinity.
\end{proposition}

\begin{proof}
Each one's only neighbour with a route to C is the other, so with $c$ the link cost, $a_{r} = b_{r-1} + c$ and $b_r =
a_{r-1} + c$ (capped at infinity). Hence $\max(a_r, b_r) = \max(a_{r-1}, b_{r-1}) + c$.
\end{proof}

\begin{omfigure}
\begin{tikzpicture}
\begin{axis}[omrecord, width=0.62\linewidth, height=4.6cm, xlabel={round after the failure},
    ylabel={distance to C}, ymin=0, ymax=17.5, xmin=0, xmax=16, legend pos=south east,
    legend style={font=\scriptsize}]
  \addplot[const plot, omDef, thick, mark=*, mark size=1.2pt] table[x=round, y=A]
    {code/bachelor-2/18-network-layer/figures/count.dat};
  \addlegendentry{A}
  \addplot[const plot, omThm, thick, mark=square*, mark size=1.2pt] table[x=round, y=B]
    {code/bachelor-2/18-network-layer/figures/count.dat};
  \addlegendentry{B}
  \draw[dashed, gray] (axis cs:0,16) -- (axis cs:16,16) node[pos=0.12, below, font=\scriptsize] {infinity: 16};
\end{axis}
\end{tikzpicture}
\par\smallskip{\small A and B take turns raising their distance to C by 2, each believing the other still has a
route. 15 rounds pass before both declare C unreachable.}
\end{omfigure}

% ledger: b2:rip-infinity
That is why distance-vector protocols take a small number for infinity: RIP, the classic one, uses 16 and sends its
tables every 30 seconds, so a count to infinity could take minutes. A remedy, \emph{split horizon with poisoned
reverse}, has a \omterm{def:b2:network-layer:routing}{router} tell the neighbour it routes through that its distance is infinite; it cures this two-router
loop, but not every loop (\cref{exo:b2:network-layer:10}).

\subsection{Link state}

% ledger: b2:ospf-link-state
\begin{definition}[Link-state routing]\label{def:b2:network-layer:ls}
In \emph{link-state routing}\index{link state routing}, each \omterm{def:b2:network-layer:routing}{router} floods to every other \omterm{def:b2:network-layer:routing}{router} a description of its
own links, so that every \omterm{def:b2:network-layer:routing}{router} holds the same map of the whole network, and each computes its shortest paths from
itself with Dijkstra's algorithm.
\end{definition}

 When a link fails, the two \omterm{def:b2:network-layer:routing}{routers} at its ends flood the news, and every \omterm{def:b2:network-layer:routing}{router}
computes again, once. OSPF, the common link-state protocol, works this way. The cost is that every \omterm{def:b2:network-layer:routing}{router} stores the
whole map and runs Dijkstra on each change; the gain is fast, loop-free convergence.

\omcode{code/bachelor-2/18-network-layer/listings/link_state.py}{5}{19}{Each \omterm{def:b2:network-layer:routing}{router}'s computation in link state:
Dijkstra, keeping the first hop}

\subsection{Between operators}

% ledger: b2:bgp-as
The Internet is a network of networks.

\begin{definition}[Autonomous system, path-vector routing]\label{def:b2:network-layer:as}
An \emph{autonomous system}\index{autonomous system} is a set of \omterm{def:b2:network-layer:routing}{routers} under one technical administration, which
uses its own \omterm{def:b2:network-layer:routing}{routing} protocol inside. In \emph{path-vector routing}\index{path-vector routing}, each announcement of a
\omterm{def:b2:network-layer:prefix}{prefix} carries the list of autonomous systems it passed through, so that a system that finds itself in the list knows
there is a loop and ignores the announcement.
\end{definition}

Between \omterm{def:b2:network-layer:as}{autonomous systems} runs BGP, a path-vector protocol. The choice among
routes is a matter of policy (which neighbour is a customer, which one is paid) rather than of shortest paths.

\omimg{book4/interface-message-processor.jpg}{0.4\linewidth}{An Interface Message Processor, one of the first packet
switches: in the network that grew into the Internet, each of these minicomputers forwarded packets between its
host computers and its neighbours. Photograph: ArnoldReinhold, CC BY 4.0.}

\section{Exercises}\label{sec:b2:network-layer:exercises}

\begin{exercise}[$\star$]\label{exo:b2:network-layer:1}
Give the mask, network address, broadcast address and number of host addresses of \code{198.51.100.77/27} and of
\code{192.0.2.130/25}.
\end{exercise}

\begin{exercise}[$\star$]\label{exo:b2:network-layer:2}
This mask is one bit short. What \omterm{def:b2:network-layer:prefix}{prefix} does it really compute, and is the answer to the membership test right?
\omcode{code/bachelor-2/18-network-layer/listings/exo_mask.py}{5}{6}{A shifted mask}
\omoutput{code/bachelor-2/18-network-layer/outputs/exo_mask.txt}{For 203.0.113.77/26.}
\end{exercise}

\begin{exercise}[$\star$]\label{exo:b2:network-layer:3}
Write \code{contains(prefix, address)} for \omterm{def:b2:network-layer:prefix}{prefixes} and addresses in text, using \code{cidr.py}.
\end{exercise}

\begin{exercise}[$\star$]\label{exo:b2:network-layer:4}
How many machines can a \code{/26} hold? Split \code{203.0.113.0/24} into the smallest \omterm{def:b2:network-layer:prefix}{prefixes} for four offices of
50, 20, 10 and 2 machines, without overlap.
\end{exercise}

\begin{exercise}[$\star\star$]\label{exo:b2:network-layer:5}
A \omterm{def:b2:network-layer:routing}{router} has the routes \code{0.0.0.0/0} $\to$ W, \code{198.51.100.0/24} $\to$ X, \code{198.51.100.64/26}
$\to$ Y and \code{198.51.100.96/27} $\to$ Z. Where do packets for \code{198.51.100.70}, \code{198.51.100.100},
\code{198.51.100.130} and \code{192.0.2.1} go?
\end{exercise}

\begin{exercise}[$\star\star$]\label{exo:b2:network-layer:6}
After the \omterm{def:b2:network-layer:ttl}{time to live} is lowered, this \omterm{def:b2:network-layer:routing}{router} computes the \omterm{def:b2:link-layer:edc}{checksum} without first clearing the field. What does the
next \omterm{def:b2:network-layer:routing}{router} see, and how does the error show in the output?
\omcode{code/bachelor-2/18-network-layer/listings/exo_hdrsum.py}{4}{7}{A \omterm{def:b2:link-layer:edc}{checksum} over itself}
\omoutput{code/bachelor-2/18-network-layer/outputs/exo_hdrsum.txt}{The receiver's sum.}
\end{exercise}

\begin{exercise}[$\star\star$]\label{exo:b2:network-layer:7}
Write \code{aggregate(prefixes)}, which replaces a set of \omterm{def:b2:network-layer:prefix}{prefixes} with the same next hop by the fewest \omterm{def:b2:network-layer:prefix}{prefixes}
covering exactly the same addresses: drop a \omterm{def:b2:network-layer:prefix}{prefix} inside another, merge two halves of a \omterm{def:b2:network-layer:prefix}{prefix}. Test it against
\code{ipaddress.collapse_addresses}.
\end{exercise}

\begin{exercise}[$\star\star$]\label{exo:b2:network-layer:8}
Prove \cref{thm:b2:network-layer:dv-converge}.
\end{exercise}

\begin{exercise}[$\star\star$]\label{exo:b2:network-layer:9}
Write \code{fragment(payload_len, mtu)}, returning the offset, size and flag of each fragment, check it against
\cref{out:b2:network-layer:fragment}, and give by hand the fragments of 3000 bytes for an \omterm{def:b2:network-layer:fragmentation}{MTU} of 1006. Why must
every fragment but the last carry a multiple of 8 bytes?
\end{exercise}

\begin{exercise}[$\star\star\star$]\label{exo:b2:network-layer:10}
Add split horizon with poisoned reverse to the rounds of distance vector. Show that it cures the line A--B--C when
B--C fails, but not a triangle A, B, C with D behind C, when C--D fails. Explain the second case.
\end{exercise}

\begin{exercise}[$\star\star\star$]\label{exo:b2:network-layer:11}
Compare the cost of a lookup in a trie, in a linear table, and by binary search in a sorted list of \omterm{def:b2:network-layer:prefix}{prefixes}, which
must also handle nested \omterm{def:b2:network-layer:prefix}{prefixes}. Relate your estimates to \cref{out:b2:network-layer:bench}.
\end{exercise}

\begin{exercise}[$\star\star\star$]\label{exo:b2:network-layer:12}
Write \code{compress(groups)}, which gives the recommended text form of an \omterm{def:b2:network-layer:ipv6}{IPv6 address} from its eight groups, and
test it against \code{ipaddress} on many addresses with runs of zeros.
\end{exercise}

\section{Problem: Routing the town's network}\label{sec:b2:network-layer:problem}

\begin{problem}[{Weekend problem --- Routing the town's network}]\label{pb:b2:network-layer:1}
The town puts a \omterm{def:b2:network-layer:routing}{router} at each of its 30 places, joined along its 53 roads, and numbers them from the block
\code{198.51.100.0/24}.

\needspace{4\baselineskip}\medskip\noindent\textbf{Part I --- An addressing plan.}\par\nopagebreak
\begin{enumerate}
\item Each \omterm{def:b2:network-layer:routing}{router} gets one address, \code{198.51.100.}$k$ for place $k$. How many addresses of the block are left?
\item Each place also has a small network of at most 5 machines. What is the smallest \omterm{def:b2:network-layer:prefix}{prefix} for one, and how many
such \omterm{def:b2:network-layer:prefix}{prefixes} fit in a \code{/24}?
\item Could the 30 small networks all fit in \code{198.51.100.0/24} besides the \omterm{def:b2:network-layer:routing}{routers}' addresses? Justify.
\item The town gets a second block, \code{203.0.113.0/24}, and gives place $k$ the \code{/29} that starts at
\code{203.0.113.}$8k$. Check that the 30 \omterm{def:b2:network-layer:prefix}{prefixes} fit, and give the \omterm{def:b2:network-layer:prefix}{prefix} of place 12.
\item Write the network and broadcast addresses of that \omterm{def:b2:network-layer:prefix}{prefix}.
\item Why does this plan make \omterm{def:b2:network-layer:lpm}{forwarding tables} easy to read?
\end{enumerate}

\needspace{4\baselineskip}\medskip\noindent\textbf{Part II --- Forwarding tables.}\par\nopagebreak
\begin{enumerate}[resume, beginpenalty=10000]
\item In \cref{out:b2:network-layer:hook}, which \omterm{def:b2:network-layer:routing}{router} forwards the packet to the Town Square \omterm{def:b2:network-layer:routing}{router}?
\item How many entries does a \omterm{def:b2:network-layer:routing}{router}'s table need with one \code{/29} per place, and how many if the town announces
itself to the outside as one \omterm{def:b2:network-layer:prefix}{prefix}?
\item Write the two entries a \omterm{def:b2:network-layer:routing}{router} at the edge of town needs for the outside world.
\item A packet for \code{203.0.113.97} arrives. Which place's network is it in?
\item Count the steps of a trie lookup for that address in a table of the 30 \code{/29}s and a default route.
\item Why is the \omterm{def:b2:network-layer:ttl}{time to live} set to 64 rather than to the 9 hops the hook needed?
\end{enumerate}

\needspace{4\baselineskip}\medskip\noindent\textbf{Part III --- Distance vector against link state.}\par\nopagebreak
\begin{enumerate}[resume, beginpenalty=10000]
\item With road lengths as costs, distance vector stops changing after 8 rounds. What does
\cref{thm:b2:network-layer:dv-converge} say the number of rounds is at most, and what does 8 measure?
\item With hop counts as costs, it takes 7 rounds. Why fewer?
\item How many messages does one round of distance vector send in the town's network?
\item How many link descriptions does each \omterm{def:b2:network-layer:routing}{router} store in link state?
\item Estimate the work of one Dijkstra run on 30 \omterm{def:b2:network-layer:routing}{routers} and 53 links.
\item Which design would you choose for the town, and why?
\end{enumerate}

\needspace{4\baselineskip}\medskip\noindent\textbf{Part IV --- A router fails.}\par\nopagebreak
\begin{enumerate}[resume, beginpenalty=10000]
\item The \omterm{def:b2:network-layer:routing}{router} at the Lighthouse Café breaks down. Why do the others not learn at once that it is unreachable?
\omoutput[1-1,4-9]{code/bachelor-2/18-network-layer/outputs/town_network.txt}{After the failure, with hop counts.}
\item Explain the first three lines: why does the nearest \omterm{def:b2:network-layer:routing}{router}'s distance go 1, 2, 3, 4?
\item Why does it take 15 rounds, whatever the \omterm{def:b2:network-layer:routing}{router}'s distance before?
\item With RIP's 30 seconds per round, how long is the café's network wrongly announced as reachable?
\item What does link state do instead, and how long does it take?
\item State the result: count to infinity, 15 rounds to reach the limit 16, against one recomputation in link state.
\end{enumerate}
\end{problem}
